
\documentclass{article}
\usepackage{colortbl}
\usepackage{makecell}
\usepackage{multirow}
\usepackage{supertabular}

\begin{document}

\newcounter{utterance}

\centering \large Interaction Transcript for game `eqbench', experiment `eq\_bench', episode 155 with reflexion{-}llama{-}3.1{-}8b{-}instruct.
\vspace{24pt}

{ \footnotesize  \setcounter{utterance}{1}
\setlength{\tabcolsep}{0pt}
\begin{supertabular}{c@{$\;$}|p{.15\linewidth}@{}p{.15\linewidth}p{.15\linewidth}p{.15\linewidth}p{.15\linewidth}p{.15\linewidth}}
    \# & \multicolumn{2}{c}{Player} && \multicolumn{2}{c}{Game Master} \\
    \hline

    \theutterance \stepcounter{utterance}  
    & & & \multicolumn{4}{p{0.6\linewidth}}{
        \cellcolor[rgb]{0.9,0.9,0.9}{
            \makecell[{{p{\linewidth}}}]{
                \texttt{\tiny{[P1$\langle$GM]}}
                \texttt{You have attempted this same task before and did not succeed. These are the plans you wrote after those attempts.} \\
\texttt{Trial \#1: New plan: I failed to accurately predict Rev. Montague's emotional state at the end of the dialogue. Upon re{-}reading the conversation, I noticed that Dr. LeGuin's final statement, "It's found in evidence, in reason," seems to be a decisive blow to Rev. Montague's argument. Rev. Montague's previous statements, such as "It's a guiding principle, a moral compass" and "Faith doesn't demand blindness, Iris. It demands trust, surrender," suggest that he is deeply invested in his faith and views it as a source of strength and guidance. Dr. LeGuin's rejection of faith as a guiding principle and her emphasis on evidence and reason would likely leave Rev. Montague feeling dismissed, as if his deeply held beliefs are being dismissed or invalidated. I should have scored Dismissed higher, at 9 or 10, and Resentful lower, at 4 or 5, as Rev. Montague's tone and language suggest a sense of frustration and defensiveness, but not necessarily resentment.} \\
\\ 
\\ 
\\ 
\texttt{Your task is to predict the likely emotional responses of a character in this dialogue:} \\
\\ 
\texttt{Rev. Montague: You see, Iris, faith is not a choice, but a divine gift. It's a beacon of light in a world often shrouded in darkness.} \\
\texttt{Dr. LeGuin: A beacon of light, Luther? Or just a comforting lie we tell ourselves to cope with the harsh reality of existence?} \\
\texttt{Rev. Montague: It's more than that. It's a guiding principle, a moral compass. It's what makes us human, Iris.} \\
\texttt{Dr. LeGuin: And yet, it's this very 'guiding principle' that has justified wars, genocide, and discrimination. How human is that, Luther?} \\
\texttt{Rev. Montague: Those are the failures of men, not faith. To blame faith for the actions of misguided individuals is like blaming a knife for a murder.} \\
\texttt{Dr. LeGuin: A knife doesn't claim to be a moral compass, Luther. And it doesn't demand blind obedience.} \\
\texttt{Rev. Montague: Faith doesn't demand blindness, Iris. It demands trust, surrender. It's a relationship, not a dictatorship.} \\
\texttt{Dr. LeGuin: A relationship with a non{-}existent entity. How convenient. } \\
\texttt{Rev. Montague: Just because you can't see something, doesn't mean it doesn't exist. Can you see your thoughts, Iris?} \\
\texttt{Dr. LeGuin: That's a false equivalence, Luther. Thoughts can be measured, analyzed. Your God, on the other hand...} \\
\texttt{Rev. Montague: Is beyond human comprehension. That's the point, Iris. We're not meant to understand everything.} \\
\texttt{Dr. LeGuin: And there lies the crux of our disagreement, Luther. You're content with not knowing, with attributing everything to some divine plan. I, on the other hand, am not. } \\
\texttt{Rev. Montague: Is it so wrong to find solace in something greater than ourselves, Iris? Is it so wrong to believe in hope, in love, in redemption?} \\
\texttt{Dr. LeGuin: It's not wrong, Luther. It's just... futile. } \\
\texttt{Rev. Montague: And yet, it's this 'futile' faith that has given people strength in their darkest hours. Can your logic do that, Iris?} \\
\texttt{Dr. LeGuin: My logic doesn't offer false hope, Luther. It offers clarity, understanding. It offers truth.} \\
\texttt{Rev. Montague: And what is truth, Iris? Your truth? My truth? Or some universal truth that we're all blindly searching for?} \\
\texttt{Dr. LeGuin: There's no 'your truth' or 'my truth', Luther. There's just the truth. And it's not found in ancient texts or divine revelations. It's found in evidence, in reason.} \\
\texttt{{[}End dialogue{]}} \\
\\ 
\texttt{At the end of this dialogue, At this point, Rev. Montague would feel...} \\
\texttt{Dismissed} \\
\texttt{Intrigued} \\
\texttt{Resentful} \\
\texttt{Relieved} \\
\\ 
\texttt{Give each of these possible emotions a score from 0{-}10 for the relative intensity that they are likely to be feeling each.} \\
\\ 
\texttt{You must output in the following format, including headings (of course, you should give your own scores), with no additional commentary:} \\
\\ 
\texttt{Dismissed: <score>} \\
\texttt{Intrigued: <score>} \\
\texttt{Resentful: <score>} \\
\texttt{Relieved: <score>} \\
\\ 
\\ 
\texttt{{[}End of answer{]}} \\
\\ 
\texttt{Remember: zero is a valid score, meaning they are likely not feeling that emotion. You must score at least one emotion > 0.} \\
\\ 
\texttt{Your answer:} \\
            }
        }
    }
     \\ \\

    \theutterance \stepcounter{utterance}  
    & \multicolumn{4}{p{0.6\linewidth}}{
        \cellcolor[rgb]{0.561,0.737,0.561}{
            \makecell[{{p{\linewidth}}}]{
                \texttt{\tiny{[P1$\rangle$GM]}}
                \texttt{Dismissed: 9} \\
\texttt{Intrigued: 0} \\
\texttt{Resentful: 2} \\
\texttt{Relieved: 0} \\
            }
        }
    }
    & & \\ \\

    \theutterance \stepcounter{utterance}  
    & & & \multicolumn{2}{p{0.3\linewidth}}{
        \cellcolor[rgb]{0.9,0.9,0.9}{
            \makecell[{{p{\linewidth}}}]{
                \texttt{\tiny{[GM$|$GM]}}
                \texttt{Dismissed: 9} \\
\texttt{Intrigued: 0} \\
\texttt{Resentful: 2} \\
\texttt{Relieved: 0} \\
            }
        }
    }
    & & \\ \\

    \theutterance \stepcounter{utterance}  
    & & & \multicolumn{2}{p{0.3\linewidth}}{
        \cellcolor[rgb]{0.9,0.9,0.9}{
            \makecell[{{p{\linewidth}}}]{
                \texttt{\tiny{[GM$|$GM]}}
                \texttt{Dismissed: 7} \\
\texttt{Intrigued: 0} \\
\texttt{Resentful: 5} \\
\texttt{Relieved: 0} \\
            }
        }
    }
    & & \\ \\

    \theutterance \stepcounter{utterance}  
    & & & \multicolumn{2}{p{0.3\linewidth}}{
        \cellcolor[rgb]{0.9,0.9,0.9}{
            \makecell[{{p{\linewidth}}}]{
                \texttt{\tiny{[GM$|$GM]}}
                \texttt{game\_result = LOSE} \\
            }
        }
    }
    & & \\ \\

    \theutterance \stepcounter{utterance}  
    & & \multicolumn{4}{p{0.6\linewidth}}{
        \cellcolor[rgb]{0.95,0.95,0.95}{
            \makecell[{{p{\linewidth}}}]{
                \texttt{\tiny{[P1$\leftrightarrow$P]}}
                \texttt{Post{-}episode Reflexion (not sent to GM):} \\
\texttt{New plan: I failed to accurately predict Rev. Montague's emotional state at the end of the dialogue. Upon re{-}reading the conversation, I noticed that Dr. LeGuin's final statement, "It's found in evidence, in reason," seems to be a decisive blow to Rev. Montague's argument. Rev. Montague's previous statements suggest that he is deeply invested in his faith and views it as a source of strength and guidance. Dr. LeGuin's rejection of faith as a guiding principle and her emphasis on evidence and reason would likely leave Rev. Montague feeling dismissed, as if his deeply held beliefs are being dismissed or invalidated.} \\
\\ 
\texttt{To improve my prediction, I will pay closer attention to the tone and language used by Rev. Montague throughout the dialogue, particularly in his responses to Dr. LeGuin's challenges to his faith. I will also consider the implications of Dr. LeGuin's final statement and how it might affect Rev. Montague's emotional state.} \\
\\ 
\texttt{In terms of specific actions, I will:} \\
\\ 
\texttt{1. Analyze Rev. Montague's language and tone to determine the level of his emotional investment in his faith.} \\
\texttt{2. Consider the potential impact of Dr. LeGuin's final statement on Rev. Montague's emotional state.} \\
\texttt{3. Use this information to adjust my scores for the possible emotions, with a focus on accurately capturing Rev. Montague's likely emotional state at the end of the dialogue.} \\
\\ 
\texttt{By taking these steps, I hope to improve my prediction and provide a more accurate assessment of Rev. Montague's emotional state.} \\
            }
        }
    }
    & \\ \\

\end{supertabular}
}

\end{document}
